\begin{explanationbox}
	\n\textbf{\textcolor{CExplanation}{一句话直觉}}\n\n海伦公式用半周长将三边整合成一个对称代数式，直接算出三角形面积，免去作高或查角。\n\n

	\medskip
	\n\textbf{\textcolor{CExplanation}{核心拆解}}\n
	\begin{description}[style=nextline, leftmargin=2.8em, labelsep=0.5em, itemsep=0.45em]
		\n
		\item[\textbf{要点一（半周长远算）：}]
			$p = \frac{a+b+c}{2}$ 是三边之和的一半，对称地包含所有边长。\n
		\item[\textbf{要点二（根号乘积）：}]
			$\sqrt{p(p-a)(p-b)(p-c)}$ 中每个因子都是非负数，保证了结果为有效面积。\n
	\end{description}
	\n\n

	\medskip
	\n\textbf{\textcolor{CExplanation}{几何本质（边长代数化）}}\n\n面积原本是底乘高等几何操作，海伦公式将其转化为关于边长的对称四次多项式开平方，本质上是三角形几何性质在代数中的投影。\n\n

	\medskip
	\n\textbf{\textcolor{CExplanation}{代数意义（无几何辅助）}}\n\n海伦公式给出了纯代数途径：只需已知三边数值，即可直接计算面积，无需测量高或角度，非常适合解析几何与代数综合题。\n\n

	\medskip
	\n\textbf{\textcolor{CExplanation}{考点价值（解题工具）}}\n
	\begin{itemize}[leftmargin=*, itemsep=0.5em, topsep=0.4em]
		\n
		\item \textbf{考法一（直接套公式）：} 已知三边数值，直接代入海伦公式求面积。\n
		\item \textbf{考法二（间接求解）：} 已知三边关系（如比例或方程），先求边长再计算面积。\n
	\end{itemize}
	\n\n

	\medskip
	\n\textbf{\textcolor{CExplanation}{顿悟点}}\n\n海伦公式中 $p-a = \frac{b+c-a}{2}$，与三角形不等式直接对应：若某边大于另两边和，$p-a$ 为负，根号无意义；正因如此，公式自带了“三角形存在条件”。\n\n

	\medskip
	\n\textbf{\textcolor{CExplanation}{使用场景}}\n
	\begin{itemize}[leftmargin=*, itemsep=0.5em, topsep=0.4em]
		\n
		\item \textbf{场景一（直接已知三边）：} 给出三边长度，用海伦公式求面积。\n
		\item \textbf{场景二（三边含参数）：} 三边用字母表示时，可先求半周长再代入，化简代数式。\n
	\end{itemize}
	\n
\end{explanationbox}
